%! TeX root: ../article.tex
\section{Resultados y discusión}
\label{sec:results}

\setlength{\emergencystretch}{6em}
\sloppy
\setlength{\tabcolsep}{4pt}

\subsection{Catálogo de experimentos}
\label{subsec:catalogo}
Antes de discutir los resultados, la tabla~\ref{tab:catalogo-experimentos}
resume los experimentos ejecutados. Cada uno varía un único eje alrededor de
una configuración base común para que las comparaciones sean justas; el eje
variado, los componentes fijos y los tamaños de banco $N$ evaluados se recogen
en la tabla. Las subsecciones siguientes se refieren a cada barrido por su
identificador.

{\footnotesize
\begin{longtable}{|p{0.6cm}|p{3.9cm}|p{3.6cm}|p{4.7cm}|} \hline
  \textbf{Id} & \textbf{Elementos variados} & \textbf{Fijo} & \textbf{$N$ / repeticiones} \\ \hline
  \endhead
  E1 & Extractor & UMAP-10, HDBSCAN \texttt{mcs}=5, identidad, SQLite & $N\in\{250, 1000, 6416\}$, $K{=}3$ \\ \hline
  E2 & Reducción & Extractor base, HDBSCAN \texttt{mcs}=5 & $N\in\{250, 1000, 6416\}$, $K{=}3$ \\ \hline
  E3 & Agrupamiento & Extractor base, UMAP-10 & $N\in\{250, 1000, 6416\}$, $K{=}3$ \\ \hline
  E4 & HDBSCAN \texttt{min\_cluster\_size} & Extractor base, UMAP-10 & $N\in\{250, 1000, 6416\}$, $K{=}3$ \\ \hline
  E5 & Segmentador & \textit{Pipeline} base & $N\in\{250, 1000, 6416\}$, $K{=}3$ \\ \hline
  E6 & Almacenamiento & \textit{Pipeline} base, búsqueda activa & $N\in\{500, 1000, 2500, 6416\}$, $K{=}5$ \\ \hline
  E7 & $N$ $\times$ agrupamiento & Extractor base, UMAP-10, SQLite & $N\in\{250, 500, 1000, 2000,$ $3500, 6416\}$, $K{=}5$ \\ \hline
  E8a & Extractor $\times$ reducción $\times$ agrupamiento sobre etiquetas de estilo & SQLite, identidad, sin límite & 294 recortes, 4 clases, $K{=}3$ \\ \hline
  E8b & Extractor $\times$ reducción $\times$ agrupamiento sobre etiquetas de autoría & SQLite, identidad, sin límite & 185 recortes, 87 autores, $K{=}3$ \\ \hline
  E9 & Parámetros HNSW & \textit{Pipeline} base, PostgreSQL & $N{=}6416$, $K{=}5$ \\ \hline
  \caption[Catálogo de barridos experimentales]{Catálogo de barridos experimentales. El detalle metodológico común
  (métricas, cronometría, repeticiones) se describe en \ref{sec:metodologia}.}
  \label{tab:catalogo-experimentos}
\end{longtable}
}

Los experimentos se ejecutaron en un portátil con CPU AMD Ryzen 7 5800H, GPU
NVIDIA GeForce GTX 1650 y 16 GB de RAM, sobre Python 3.14 gestionado con
\texttt{uv} y, como dependencias nucleares, PyTorch (fijado a 2.10.0),
\textit{open-clip-torch}, \textit{ultralytics}, \textit{scikit-learn},
\textit{umap-learn} y \textit{SQLAlchemy} con \textit{pgvector}. Las versiones
completas de las dependencias se recogen en el anexo~\ref{anexo:repro}
(tabla~\ref{tab:dependencias}).

\subsection{Calibración del \textit{pipeline}}
\label{subsec:res-calibracion}
En esta sección se recorren las decisiones tomadas para el diseño de lo que se
considera el \textit{pipeline} base del proyecto, la configuración sobre la que
se asientan los estudios de coste y calidad supervisada en las secciones
siguientes. La elección de reducción y agrupamiento queda decidida por la triada
de barridos E2, E3 y E4, cuyo desarrollo completo se recoge en el
anexo~\ref{anexo:resultados} y cuyas conclusiones se resumen aquí
(\ref{subsubsec:res-reduccion-resumen}); la del segmentador se discute a
continuación. La excepción es la elección del extractor, que se basa en un
criterio de calidad supervisada (E8a) que se discute en la sección de
resultados supervisados (\ref{subsec:res-supervisada}).

La \textit{calidad intrínseca} (silueta~$\uparrow$, CH~$\uparrow$,
DB~$\downarrow$; definidas en \ref{subsubsec:eval-intrinseca}) se lee siempre
de forma conjunta con el número de clústeres, el ruido y el coeficiente de
variación de tamaños. Leer la silueta o el CH de forma aislada premia
configuraciones degeneradas (pocos clústeres grandes y separados, o particiones
binarias triviales) y produce inversiones consistentes del ranking.

\subsubsection{Extractor de características}
\label{subsubsec:res-extractor}
El primer paso del pipeline es la extracción de características. Para esto se
prueban doce extractores: cuatro CNN preentrenadas en ImageNet
(\texttt{resnet50}, \texttt{vgg16}, \texttt{inception\_v3},
\texttt{mobilenet\_v3}), \texttt{dinov2\_vits14}, \texttt{clip\_vit\_b32}, dos
\textit{backbones} YOLO (\texttt{yolon}, \texttt{yolom}) y la matriz completa
2$\times$2 de cabezas \textit{fine-tuned} (autor y estilo $\times$ DINOv2 y
MobileNetV3); la tabla~\ref{tab:e1-quality-6416} recoge la calidad intrínseca a
$N=6416$.

{\footnotesize
\begin{longtable}{|p{5.4cm}|p{1.0cm}|p{1.0cm}|p{1.0cm}|p{1.0cm}|p{1.0cm}|p{1.0cm}|} \hline
  \textbf{Extractor} & \textbf{cls} & \textbf{ruido} & \textbf{sil} & \textbf{CH} & \textbf{DB} & \textbf{csCV} \\ \hline
  \endhead
  \texttt{mobilenet\_v3}                  & 820 & \textbf{0.078} & \textbf{0.260} & 13.1 & 1.48 & \textbf{0.52} \\ \hline
  \texttt{resnet50}                       & 794 & 0.086 & 0.257 & 15.6 & \textbf{1.46} & 0.71 \\ \hline
  \texttt{vgg16}                          & 773 & 0.095 & 0.228 & 13.4 & 1.55 & 0.54 \\ \hline
  \texttt{mob.\_graffiti\_author\_head}   & 779 & 0.124 & 0.227 & 13.4 & 1.53 & 0.53 \\ \hline
  \texttt{inception\_v3}                  & 724 & 0.118 & 0.213 & 11.5 & 1.61 & 0.64 \\ \hline
  \texttt{dinov2\_graffiti\_author\_head} & 631 & 0.143 & 0.187 & 17.9 & 1.59 & 0.81 \\ \hline
  \texttt{dinov2\_vits14}                 & 632 & 0.123 & 0.185 & 19.2 & 1.58 & 0.92 \\ \hline
  \texttt{clip\_vit\_b32}                 & 734 & 0.150 & 0.170 & 11.9 & 1.66 & 0.74 \\ \hline
  \texttt{mob.\_graffiti\_style\_head}    & 629 & 0.214 & 0.164 & 19.7 & 1.74 & 0.58 \\ \hline
  \texttt{dinov2\_graffiti\_style\_head}  & 562 & 0.202 & 0.155 & 27.9 & 1.70 & 0.74 \\ \hline
  \texttt{yolom}                          & 480 & 0.253 & 0.131 & 32.2 & 1.75 & 0.74 \\ \hline
  \texttt{yolon}                          & 348 & 0.341 & 0.085 & \textbf{51.2} & 1.91 & 0.81 \\ \hline
  \caption[E1, calidad intrínseca a $N=6416$]{E1, calidad intrínseca a $N=6416$. Las CNN entrenadas en ImageNet encabezan la silueta; el extractor base
  (\texttt{dinov2\_graffiti\_style\_head}) presenta un CH (27.9) solo superado
  por los \textit{backbones} YOLO, pese a una silueta inferior a la de su
  \textit{backbone} DINOv2. \glosacalidad}
  \label{tab:e1-quality-6416}
\end{longtable}
}

El resultado más llamativo del barrido es la inversión entre silueta y
Calinski--Harabasz sobre ambas cabezas estilísticas: la silueta cae respecto a
su \textit{backbone} (DINOv2: $0.185\to0.155$; MobileNetV3: $0.260\to0.164$)
mientras que el CH sube (DINOv2: $19.2\to27.9$; MobileNetV3: $13.1\to19.7$). Es
coherente con el entrenamiento contrastivo, que optimiza la separación entre
clústeres (CH) pero no la cohesión local de cada punto (silueta), además de ser
agravado por un cambio de dominio: las cabezas se entrenaron con recortes, no
con las imágenes completas evaluadas aquí.

Las CNN entrenadas en ImageNet lideran la silueta (a $N=6416$:
\texttt{mobilenet\_v3}=0.260, \texttt{resnet50}=0.257), con ruido bajo y
clústeres balanceados, pero esto no implica que sean mejores agrupando grafiti:
es probable que estén atrupando por apariencia genérica (textura, color,
composición) sin codificar información del dominio. Los modelos de YOLO, en
cambio, pierden en todos los ejes salvo el CH, confirmando que sus
características de detección no transfieren al agrupamiento por apariencia.

En coste (tabla~\ref{tab:e1-cost-6416}, anexo~\ref{anexo:detalle}),
\texttt{mobilenet\_v3} es el más barato en velocidad y VRAM (7.93~ips, 32~MB) y
la cabeza añade un sobrecoste pequeño ($\sim$8\,\%); los \textit{backbones}
YOLO son los más caros ($\sim$825~MB de RSS) y VGG16/CLIP los más exigentes en
VRAM (552 y 592~MB).

\subsubsection{Reducción y agrupamiento base}
\label{subsubsec:res-reduccion-resumen}
Las dos etapas intermedias del \textit{pipeline} base se fijan con tres
barridos cuyo desarrollo completo (tablas, figuras y discusión) se describen en
el anexo~\ref{anexo:resultados}; aquí solo se recogen sus conclusiones. El
barrido de \textbf{reducción} (E2) señala a \textbf{UMAP con
\texttt{n\_components=10}} como la única familia que produce a la vez un número
de clústeres no trivial, ruido moderado y tamaños balanceados: la identidad y
PCA-50 obtienen mejor silueta aislada, pero descartando dos tercios del banco
como ruido, y la identidad además multiplica por $\sim$20 el coste de
agrupamiento. UMAP-50 iguala a UMAP-10 en calidad al doble de coste.

El barrido de \textbf{agrupamiento} (E3) sitúa a \textbf{HDBSCAN} como el único
algoritmo que combina silueta competitiva, escalado sensato de $k$ con $N$,
ruido moderado, tamaños balanceados y coste despreciable. OPTICS lo supera
ligeramente en calidad pero a $\sim$25 veces el tiempo y $\sim$2500 veces la
memoria a $N=6416$. Por último, el barrido de \textbf{granularidad} (E4) fija
\textbf{\texttt{min\_cluster\_size}=5}: los valores altos maximizan silueta y
CH pero colapsan a dos clústeres degenerados, y la elección por un grano más
fino se ve reforzada después por las métricas supervisadas de E8a.

\subsubsection{Impacto del segmentador}
\label{subsubsec:res-segmentador}
En principio, el recorte de la imagen completa en unidades más pequeñas que
contengan únicamente un grafiti debería facilitar el agrupamiento, al eliminar
ruido de fondo y permitir que el extractor de características se centre en la
información relevante. Este cambio debería ser especialmente beneficioso para
los extractores entrenados con recortes (las cabezas de DINOv2 y MobileNetV3).
El barrido E5 compara cuatro opciones de segmentación (imagen completa,
YOLO11s, YOLO11m y YOLO11m \textit{fine-tuned}) manteniendo fijo el resto del
\textit{pipeline}.

{\footnotesize
\begin{longtable}{|p{2.5cm}|p{1.5cm}|p{1.0cm}|p{1.0cm}|p{1.1cm}|p{1.0cm}|p{1.0cm}|p{1.4cm}|} \hline
  \textbf{Segmentador} & \textbf{recortes} & \textbf{r/img} & \textbf{cls} & \textbf{ruido} & \textbf{sil} & \textbf{CH} & \textbf{ingesta (s)} \\ \hline
  \endhead
  identidad           & 6416  & 1.00 &  563 & \textbf{0.199} & \textbf{0.151} & 33.1 & \textbf{1235.5} \\ \hline
  YOLO11s             & 6590  & 1.03 &  502 & 0.225 & 0.127 & 30.4 & 1777.0 \\ \hline
  YOLO11m             & 6603  & 1.03 &  506 & 0.230 & 0.131 & 29.5 & 1811.6 \\ \hline
  YOLO11m \textit{ft} & 15362 & 2.39 & 1049 & 0.298 & 0.122 & \textbf{36.2} & 2152.3 \\ \hline
  \caption[E5, calidad y coste de ingesta a $N=6416$]{E5, calidad intrínseca y coste de ingesta a $N=6416$ (la columna
  \enquote{recortes} son los puntos efectivamente agrupados, no las imágenes leídas).
  Cada segmentador alimenta a HDBSCAN un conjunto distinto de puntos.
  Columnas: r/img=recortes por imagen, cls=número de clústeres,
  ruido=fracción de ruido~$\downarrow$, sil=silueta~$\uparrow$,
  CH=Calinski--Harabász~$\uparrow$, ingesta=tiempo de ingesta.}
  \label{tab:e5-segmenter-6416}
\end{longtable}
}

A escala completa, la identidad gana en todos los ejes que importan: silueta
más alta (0.151), ruido más bajo (0.199) y mejor coste de ingesta (1235~s vs
1777--2152~s). Los detectores base (YOLO11s y YOLO11m) recortan $\sim$1.03
cajas por imagen, prácticamente reproduciendo el conjunto original con una
ganancia marginal o negativa en calidad y un sobrecoste de ingesta del 44\%.

El detector \textit{fine-tuned} es un caso interesante: a $N=1000$ obtiene la
mejor silueta, mejor CH y mejor DB del panel, pero este liderazgo desaparece
con $N=6416$, donde obtiene casi 2.4 recortes sobre cada imagen original y
degenera el agrupamiento, con un 29.8\% de ruido, una silueta de 0.122 y un
coeficiente de variación de tamaños de 1.06, con $1049$ clústeres y una cola
larga de clústeres pequeños. El coste de ingesta también se dispara a 2152~s,
un 74\% más que la identidad.

Las razones de esta caída de calidad se vuelven más claras al consultar las
métricas de validación del modelo de detección. En primer lugar, este
obtiene una precisión de $0.849$ y un recall de $0.741$. Esto implica que,
cuando el modelo detecta un grafiti, lo hace correctamente el 84.9\% de las
veces, pero también que no anota uno de cada cuatro grafitis presentes en la
imagen. Por otro lado, obtiene valores de $0.808$ y $0.603$ para mAP@0.5 y
mAP@0.5-0.95 respectivamente. Esto nos indica que, incluso en los casos en los
que el modelo detecta un grafiti, la caja delimitadora no se ajusta con mucha
precisión al valor real, lo que sugiere que los recortes pueden ser de calidad
variable.

Con estos datos en mente, es posible que el modelo \textit{fine-tuned} produzca
recortes de calidad muy variable, muchas veces llegando a no recortar la imagen
y devolviéndola completa o dejando atrás varios grafitis. Esta variabilidad en
los datos de entrada, junto con el hecho de que el extractor base no fue
entrenado con recortes de este tipo, podría ser la causa de la caída de calidad
y el aumento de coste a medida que el tamaño del banco crece. El
anexo~\ref{anexo:detalle} (figura~\ref{fig:e5-yolo-detections}) recoge ejemplos
de detecciones del modelo obtenidas durante el entrenamiento.

Esta conclusión cierra la configuración base del trabajo y la fija como:
\texttt{dinov2\_graffiti\_style\_head}~+~UMAP-10~+~HDBSCAN
\texttt{mcs}=5~+~segmentador identidad~+~SQLite. Es la configuración sobre la
que se asientan las subsecciones siguientes.

\subsection{Coste computacional}
\label{subsec:res-coste}
Este es uno de los ejes centrales del trabajo: caracterizar empíricamente cómo
crece el coste de cada etapa con el tamaño del banco y contrastar las
pendientes observadas con las clases de complejidad teóricas. El rango
disponible ($N\le 6416$) es estrecho para separar asintóticamente $O(n\log n)$
de $O(n^2)$, por lo que las pendientes se reportan como \emph{observadas en
este régimen} y no como exponentes asintóticos.

\subsubsection{Escalabilidad por etapa}
\label{subsubsec:res-escalabilidad}
El barrido de escalabilidad (E7) mide ingesta, reducción, agrupamiento (cinco
algoritmos) y búsqueda por similitud sobre seis tamaños de banco; la
tabla~\ref{tab:e7-cost} (anexo~\ref{anexo:detalle}) consolida los tiempos por
etapa a las tres escalas extremas con HDBSCAN como agrupador de referencia, y la
tabla~\ref{tab:e7-slopes} contrasta las pendientes log-log observadas con la
complejidad teórica de referencia. La figura~\ref{fig:e7-scaling} representa
estos tiempos frente a $N$ en escala log-log.


{\footnotesize
\begin{longtable}{|p{2.5cm}|p{3.4cm}|p{2.6cm}|p{2.6cm}|} \hline
  \textbf{Etapa / Algoritmo} & \textbf{Escala teórica} & \textbf{Pendiente $N{=}250\ldots6416$} & \textbf{Pendiente $N{=}500\ldots3500$} \\ \hline
  \endhead
  Ingesta (DINOv2)       & $O(n)$            & 0.95 & --   \\ \hline
  Reducción (UMAP)       & $\sim O(n^{1.14})$ & 1.15 & 1.53 \\ \hline
  K-medias ($k$ fijo)      & $\approx O(n)$    & 0.43 & 0.46 \\ \hline
  DBSCAN                 & $O(n\log n)$--$O(n^2)$ & 0.99 & 1.10 \\ \hline
  OPTICS                 & $O(n\log n)$--$O(n^2)$ & 0.94 & 0.98 \\ \hline
  HDBSCAN                & $O(n\log n)$ amort. & 1.53 & 1.84 \\ \hline
  Aglomerativo           & $O(n^2 \log n)$   & 1.77 & 1.76 \\ \hline
  Búsqueda (SQLite)      & $O(n^2)$ por pasada & \textbf{2.10} & 2.13 \\ \hline
  \caption[E7, pendientes observadas por etapa]{E7, pendientes observadas (regresión log-log) por etapa, comparadas
  con la complejidad teórica de referencia. La ventana $500\ldots3500$ excluye
el punto $N=6416$ de reducción (anomalía \texttt{pynndescent}, ver texto).}
  \label{tab:e7-slopes}
\end{longtable}
}

Como se puede ver en la segunda tabla, el rango reducido de $N$ provoca que
varias de las etapas tengan pendientes observadas menores que la complejidad
teórica de referencia, como puede ser el caso de DBSCAN o OPTICS. Algoritmos
como HDBSCAN o el aglomerativo, con complejidades más altas, sí que presentan
pendientes superiores a 1, dejando clara su naturaleza super-lineal. El tiempo
de ejecución de estos algoritmos fue prácticamente despreciable al compararlo
con la ingesta o la búsqueda por similitud, no llegando a superar los 5~s en
ningún caso. El único cuello de botella para el agrupamiento a esta escala es
el aglomerativo, que a $N=6416$ consume $314.1$~MB de RAM, creciendo como
$O(n^2)$ y descartándolo operativamente mucho antes de que su tiempo de pared
lo haga.

La pendiente observada de la búsqueda por similitud SQLite es prácticamente la
esperada: $\approx 2.10$ sobre el rango completo y 2.13 en la ventana limpia. A
$N=6416$ una sola pasada de 6416 consultas sobre el dataset tarda 54.4 minutos,
más que la propia ingesta (21.6 min) y unas $9000$ veces más que el
agrupamiento ($3266.6 / 0.358$~s). El \textit{throughput} cae de 68 consultas
por segundo a $N=250$ a 2 consultas por segundo a $N=6416$: una caída de
35 veces. Este resultado es la motivación cuantitativa directa para el
barrido de \textit{backend} y de índice ANN
(\ref{subsubsec:res-similaritysearch}).

La anomalía más importante del barrido se encuentra en la ingesta para el
algoritmo de reducción UMAP. Este, crece hasta $N=3500$ con una pendiente de
1.15, pero a $N=6416$ cae a 6.24~s, un 47\% más rápido sobre un 83\% más de
datos. Esta caída es reproducible y se debe a que UMAP utiliza un umbral
interno para cambiar a un cálculo aproximado de vecinos basado en
\texttt{pynndescent} al superar unos pocos miles de puntos. Esto provoca que la
pendiente completa no describa un único régimen, sino dos, y es la razón por la
que se reportan dos pendientes distintas en la tabla~\ref{tab:e7-slopes}.

\begin{figure}[h!]
  \centering
  \includegraphics[width=\linewidth]{figures/e7_scalability_loglog.pdf}
  \caption[E7, tiempo de pared vs $N$ por etapa]{E7 (figura titular): tiempo de pared vs $N$ por etapa del
  \textit{pipeline}, escala log-log. La pendiente de la búsqueda por
  similitud SQLite ($\approx 2.10$) confirma el coste $O(n^2)$ por pasada.
La caída de UMAP a $N=6416$ refleja el cambio interno a \texttt{pynndescent}.}
  \label{fig:e7-scaling}
\end{figure}

\subsubsection{Búsqueda por similitud: \textit{backend} e índice ANN}
\label{subsubsec:res-similaritysearch}
El cuello $O(n^2)$ de la búsqueda SQLite identificado en E7 motiva el barrido
E6 sobre tres \textit{backends}: SQLite (exploración lineal en Python),
pgvector exacto (escaneo secuencial dentro del SGBD con el \textit{kernel} de
coseno en C) y pgvector con índice HNSW (\texttt{m}=16,
\texttt{ef\_construction}=64). La construcción del índice HNSW se cronometra
una sola vez entre la ingesta y el bucle de búsqueda, así que su coste no se
acumula sobre el tiempo de consulta. La tabla~\ref{tab:e6-storage}
(anexo~\ref{anexo:detalle}) recoge la latencia, el \textit{throughput} y el
\textit{recall} por \textit{backend} y tamaño de banco.

SQLite mantiene un \textit{throughput} aproximadamente plano ($\sim$150~ips)
con tiempo de ejecución lineal en $N$. PostgreSQL exacto eleva el
\textit{throughput} en dos órdenes de magnitud. Por debajo de $N=1000$ la
latencia está dominada por el \textit{round-trip} de la conexión. PostgreSQL
con HNSW es el único \textit{backend} cuyo \textit{throughput} \emph{crece} con
el tamaño del banco con respecto al exacto: a $N=2500$ es 100 veces más rápido
que SQLite y a $N=6416$ es 3.2 veces más rápido que el exacto y 240 veces más
rápido que SQLite, con \texttt{R@5} = 1.000 a las cuatro escalas. La
construcción del índice cuesta 0.47~s a $N=6416$, se amortiza en una sola
consulta frente a SQLite y en varias docenas frente al exacto.

Por debajo de $N\approx2500$ las tres opciones sobre PostgreSQL son
indistinguibles en latencia y HNSW solo añade complejidad sin ganancia; por
encima, HNSW pasa al frente mientras el exacto sigue creciendo linealmente. El
dataset del trabajo está por encima de ese cruce, por lo que HNSW es la
elección por defecto recomendada para la búsqueda; SQLite queda como
\textit{backend} de desarrollo de muy pequeña escala.

Para validar que la configuración por defecto del índice no es una decisión
arbitraria, un barrido complementario E9 recorre once configuraciones en torno
a la utilizada en el barrido E6 $(\texttt{m}=16, \texttt{ef\_construction}=64,
\texttt{ef\_search}=40)$ sobre el conjunto de datos completo. Los resultados se
pueden ver en la figura~\ref{fig:similaritysearch-pareto}. Ninguna cae por
debajo de \textit{recall} 0.990, y el ancla proporciona 0.321~s por consulta a
\textit{recall} 0.995, 3.4 veces más rápida que el método exacto. Alcanzar
\textit{recall} 1.000 sin penalizar la latencia es posible subiendo \texttt{m}
a 32 o \texttt{ef\_construction} a 256; ambos son parámetros para la
construcción, preferibles a \texttt{ef\_search} en un flujo de \enquote{ingesta
única, muchas consultas} porque su sobrecoste se amortiza una sola vez. A esta
escala, parece que el uso de HNSW proporciona una mejora de latencia muy
considerable sin una pérdida apreciable de calidad, eliminando el cuello de
botella $O(n^2)$ de SQLite.

\begin{figure}[h!]
  \centering
  \includegraphics[width=\linewidth]{figures/e6_e9_pareto.pdf}
  \caption[E6+E9, latencia y \texttt{R@5} por configuración HNSW]{E6+E9: latencia por consulta (eje logarítmico) y \textit{recall}@5
      por \textit{backend} y configuración HNSW a $N=6416$, una fila por
      configuración ordenada por latencia (cada barra anota su latencia y, para
      HNSW, el \textit{recall}). La configuración del barrido E6 aparece en rojo;
      SQLite y pgvector exacto se incluyen como referencia.}
  \label{fig:similaritysearch-pareto}
\end{figure}

\subsection{Validación supervisada}
\label{subsec:res-supervisada}
Las métricas extrínsecas, aplicables cuando se dispone de etiquetas sobre los
pares, las métricas extrínsecas miden directamente la concordancia entre los
clústeres y la verdad de referencia. Esta subsección hace uso de los dos
conjuntos etiquetados manualmente por autor y estilo para evaluar la capacidad
de los extractores de características para separar ambos ejes y validar las
métricas intrínsecas empleadas en el resto del trabajo. La unidad de análisis
es el recorte (recortes pre-segmentados, segmentador identidad): 294 recortes
sobre 4 estilos para E8a y 185 recortes sobre 87 autores para E8b. Los dos
diseños no son simétricos en su exposición al desplazamiento de dominio. La
cabeza estilística se entrenó sobre un conjunto \emph{mixto} de 385 recortes
(110 Salamanca + 275 Cuenca) y se evalúa sobre 294 recortes de Salamanca
disjuntos del entrenamiento, así que E8a es un \textit{held-out} dentro del
propio dominio de entrenamiento. Las cabezas de autoría, en cambio, se
entrenaron sobre 321 recortes \emph{exclusivamente} de Cuenca y se evalúan
sobre 185 recortes de Salamanca, así que E8b mezcla generalización a autores no
vistos con un desplazamiento de dominio Cuenca$\to$Salamanca.

\subsubsection{Evaluación por estilo}
\label{subsubsec:res-estilo}
El barrido E8a evalúa la capacidad de los distintos extractores para separar
los cuatro estilos de grafiti escogidos. La tabla~\ref{tab:e8a-best} recoge la
mejor celda por extractor (E8a) sobre la malla de $2$ reducciones $\times$ $6$
agrupadores, seleccionada por el ARI por defecto (sin descartar puntos de
ruido), \emph{chance-corrected} y, por tanto, comparable entre configuraciones
con distinto número de clústeres.

{\footnotesize
\begin{longtable}{|p{5.4cm}|p{1.2cm}|p{1.0cm}|p{1.0cm}|p{1.0cm}|p{3.5cm}|} \hline
  \textbf{Extractor} & \textbf{ARI} & \textbf{NMI} & \textbf{F1} & \textbf{sil} & \textbf{Celda} \\ \hline
  \endhead
  \texttt{dinov2\_graffiti\_style\_head}      & \textbf{0.722} & \textbf{0.678} & \textbf{0.806} & \textbf{0.284} & identidad + K-medias-4   \\ \hline
  \texttt{mob.\_graffiti\_style\_head}        & 0.627 & 0.539 & 0.740 & 0.272 & identidad + K-medias-4   \\ \hline
  \texttt{clip\_vit\_b32}                     & 0.368 & 0.392 & 0.556 & 0.047 & UMAP + aglomerativo    \\ \hline
  \texttt{dinov2\_vits14}                     & 0.356 & 0.374 & 0.558 & 0.135 & identidad + K-medias-4   \\ \hline
  \texttt{mobilenet\_v3}                      & 0.356 & 0.336 & 0.550 & 0.038 & UMAP + aglomerativo    \\ \hline
  \texttt{resnet50}                           & 0.334 & 0.277 & 0.548 & 0.058 & identidad + espectral  \\ \hline
  \texttt{mob.\_graffiti\_author\_head}       & 0.318 & 0.295 & 0.517 & 0.046 & UMAP + aglomerativo    \\ \hline
  \texttt{dinov2\_graffiti\_author\_head}     & 0.302 & 0.273 & 0.519 & 0.094 & UMAP + aglomerativo    \\ \hline
  \caption[E8a, mejor celda por extractor (estilo)]{E8a, mejor celda por extractor sobre la evaluación de estilo (294
  recortes, 4 clases). El ranking sitúa las cabezas de autoría por debajo del
  \textit{backbone} en ambas familias, aunque con un margen modesto
  ($-0.054$ DINOv2, $-0.038$ MobileNetV3).}
  \label{tab:e8a-best}
\end{longtable}
}

El resultado es inequívoco: la cabeza estilística de DINOv2 + identidad +
K-medias-4 obtiene ARI = 0.722 y F1 = 0.806, presenta más del doble del ARI del
\textit{backbone} y de cualquier otro extractor sin entrenar. Va seguido por la
cabeza estilística de MobileNetV3, que, aunque no llega a doblar el ARI del
siguiente extractor, sí lo mejora sustancialmente (0.627 vs 0.368). El resto de
extractores sin entrenar (CLIP, DINOv2 y MobileNetV3) se agrupan en un segundo
bloque con ARI entre 0.33 y 0.37.

Las cabezas de autoría no solo no mejoran el ARI de estilo respecto a su
propio \textit{backbone} sino que lo empeoran (DINOv2 $0.356\to 0.302$, MobileNetV3
$0.356\to 0.318$), aunque con un margen modesto que las sitúa en los puestos
finales del panel. El \textit{fine-tuning} no produce una representación
genérica para el grafiti: concentra el eje para el que fue entrenada y deja al
otro en clara desventaja. Esta matriz de transferencia diagonal-dominante es la
evidencia más limpia de que las dos cabezas codifican información genuinamente
distinta.

La señal secundaria de recuperación supervisada (búsqueda por similitud con
$k=5$, evaluada frente a la etiqueta de la consulta) confirma el ordenamiento
y prescinde de la elección de agrupador: la cabeza estilística de DINOv2 ocupa
la primera posición en los tres índices (P@5~$=0.851$, MAP@5~$=0.909$,
MRR~$=0.922$) y el orden por P@5 reproduce, en grandes trazos, el ranking por
ARI (cabezas de estilo arriba, \textit{backbones} sin entrenar en el bloque
intermedio y cabezas de autoría en la parte baja), lo que descarta que el
liderazgo de la cabeza esté inducido por la elección de agrupador. El detalle
por extractor figura en el anexo~\ref{anexo:detalle}
(tabla~\ref{tab:e8a-retrieval}).

HDBSCAN no distingue las cuatro categorías de estilo como clústeres separados,
sino que las funde en tres grupos. Una revisión manual los identifica como
\enquote{sencillos} ($\sim$150 recortes: \textit{tags} y \textit{throw-ups}
simples), \enquote{elaborados} ($\sim$120: \textit{throw-ups} elaborados y la
mayoría de \textit{pieces}) y \enquote{characters} (solo 17: personajes de
estilo cómic o manga). El anexo~\ref{anexo:detalle} muestra su distribución 2D
(figura~\ref{fig:e8a-hdbscan}) y recortes de ejemplo
(figura~\ref{fig:e8a-clusters}).

\subsubsection{Evaluación por autor}
\label{subsubsec:res-autor}
El barrido E8b evalúa la capacidad de los extractores para separar un conjunto
de imágenes por autoría. Se parte de una hipótesis simétrica a la de E8a: una
cabeza entrenada para autoría podría superar a su propio \textit{backbone} por
bastante margen, y superar a los extractores sin entrenar. La
tabla~\ref{tab:e8b-best} muestra que esta hipótesis no es correcta. El mejor
ARI global lo obtiene CLIP ViT-B/32 (0.076), y los ocho extractores caben en
una banda 0.041--0.076 cercana al agrupamiento aleatorio (la celda ganadora de
E8a vale 0.722, $\sim$10 veces más). Dentro de la familia DINOv2 la cabeza de
autoría sí supera a su \textit{backbone} (0.065 vs.\ 0.055) y a la cabeza de
estilo (0.054), aunque a escala dominada por ruido; en la familia MobileNetV3
la cabeza apenas iguala a su \textit{backbone} (0.048 vs.\ 0.046, +0.002 ARI)
y supera a la cabeza de estilo (0.041).

{\footnotesize
\begin{longtable}{|p{5.4cm}|p{1.2cm}|p{1.0cm}|p{1.0cm}|p{1.0cm}|p{3.5cm}|} \hline
  \textbf{Extractor} & \textbf{ARI} & \textbf{NMI} & \textbf{F1} & \textbf{sil} & \textbf{Celda} \\ \hline
  \endhead
  \texttt{clip\_vit\_b32}                     & \textbf{0.076} & \textbf{0.789} & 0.090 & 0.003 & UMAP + aglo-87 \\ \hline
  \texttt{dinov2\_graffiti\_author\_head}     & 0.065 & 0.647 & \textbf{0.093} & 0.052 & UMAP + aglo-30 \\ \hline
  \texttt{dinov2\_vits14}                     & 0.055 & 0.754 & 0.079 & \textbf{0.096} & identidad + aglo-87 \\ \hline
  \texttt{dinov2\_graffiti\_style\_head}      & 0.054 & 0.628 & 0.085 & 0.037 & UMAP + aglo-30 \\ \hline
  \texttt{resnet50}                           & 0.049 & 0.778 & 0.064 & 0.028 & UMAP + aglo-87 \\ \hline
  \texttt{mob.\_graffiti\_author\_head}       & 0.048 & 0.627 & 0.077 & 0.021 & UMAP + aglo-30 \\ \hline
  \texttt{mobilenet\_v3}                      & 0.046 & 0.780 & 0.060 & 0.023 & UMAP + aglo-87 \\ \hline
  \texttt{mob.\_graffiti\_style\_head}        & 0.041 & 0.629 & 0.070 & 0.020 & UMAP + aglo-30 \\ \hline
  \caption[E8b, mejor celda por extractor (autoría)]{E8b, mejor celda por extractor sobre la evaluación de autoría
  (185 recortes, 87 autores, evaluación entre ciudades Cuenca$\to$Salamanca).
  Todos los ARI caen en una banda estrecha cercana al suelo; la simetría
  esperada con E8a no se materializa.}
  \label{tab:e8b-best}
\end{longtable}
}

E8b es un resultado negativo, pero esperado. La clasificación por estilo es
mucho más sencilla que el agrupamiento por autoría, que además sufre un
desplazamiento de dominio más severo y utiliza un conjunto de peor calidad.
Obtener resultados de autoría aceptables exigiría anotar muchos más recortes y
autores (probablemente más de los que contiene el dataset completo).

El problema de recuperación sigue la misma tendencia que el anterior,
empeorando aún más la posición de las cabezas de autoría. CLIP sigue siendo el
mejor recuperador (P@5~$=0.091$), seguido de cerca por ResNet50 (0.085),
MobileNetV3 (0.083), DINOv2 base (0.081) y las dos cabezas de autoría
(MobileNetV3 0.078, DINOv2 0.077): todos dentro de un rango estrecho. En este
caso, las cabezas entrenadas para esta tarea específica se sitúan por debajo de
los extractores sin entrenar, lo que es consistente con la caída de ARI y con
la hipótesis de que el \textit{fine-tuning} para autoría no consigue una
representación útil para esta tarea. El detalle por extractor (línea base
aleatoria $P@5 \approx 0.024$) figura en el anexo~\ref{anexo:detalle}
(tabla~\ref{tab:e8b-retrieval}).

\subsubsection{Correlación interna-externa}
\label{subsubsec:res-correlacion}
La correlación de Pearson entre silueta y ARI por celda (detalle en
\ref{anexo:correlacion}) es positiva y fuerte en estilo (E8a: $r=+0.78$,
$n=48$, $p<10^{-7}$ sobre agrupadores de $k$ fijo) pero se desvanece en autoría
(E8b), donde todo el rango de ARI vive en una banda cercana al azar y el signo
del coeficiente carece de peso práctico.

El factor decisivo es la granularidad de la verdad de referencia. La silueta
recompensa pocos clústeres grandes y cohesivos: con grano grueso (4 estilos)
eso coincide con lo que el ARI premia, pero con grano fino (87 autores) ocurre
lo contrario, pues las celdas de mayor ARI tienen $k=87$ —y por fuerza una
silueta casi nula— mientras que las de mayor silueta son las particiones
degeneradas de HDBSCAN ($k=2$ o $3$) con ARI nulo. La silueta sirve, pues, para
ordenar configuraciones a granularidad de estilo, pero no de autoría.

La consecuencia práctica es la elección del extractor base: la cabeza
estilística de DINOv2, uno de los peores modelos por métricas intrínsecas en el
experimento inicial, se escoge por su liderazgo en E8a y no por el ranking
intrínseco de E1. Esta decisión es importante, puesto que haber escogido el
extractor con mejores métricas intrínsecas a $N=6416$ habría llevado a una
configuración de peor rendimiento en la tarea que motiva el propio trabajo.

